\chapter{Boolean Analysis — Invariance — Sequence}
%%% generated by the blueprint pipeline from TCSlib.BooleanAnalysis.GaussianSpace.Invariance.Sequence
\section{Overview}
A sequence of hybrid input families gives the summed replacement bound. Equality of independent coordinate laws also determines the joint vector law.

\section{Declarations}

\begin{definition}[Selection from two coordinate copies]
\lean{BooleanAnalysis.GaussianSpace.Invariance.hybridIndex}
\leanok
For integers $n,r\ge0$ and an index $i\in\{0,\ldots,n-1\}$, the hybrid selection chooses the right copy of $i$ if $i<r$ and the left copy otherwise, as an element of two disjoint copies of the coordinate set.
\end{definition}

\begin{definition}[Hybrid family after initial replacements]
\lean{BooleanAnalysis.GaussianSpace.Invariance.hybridInput}
\leanok
\uses{BooleanAnalysis.GaussianSpace.Invariance.hybridIndex}
Let $Z_{L,i},Z_{R,i}\colon\Omega\to\mathbb R$ be two indexed families, $0\le i<n$, and let $r\ge0$ be an integer. The hybrid family $H^{(r)}$ has $H^{(r)}_i=Z_{R,i}$ for $i<r$ and $H^{(r)}_i=Z_{L,i}$ for $i\ge r$.
\end{definition}

\begin{lemma}[Distinct coordinates remain distinct in hybrids]
\lean{BooleanAnalysis.GaussianSpace.Invariance.hybridIndex_injective}
\leanok
\uses{BooleanAnalysis.GaussianSpace.Invariance.hybridIndex}
\difficulty{4}
For integers $n,r\ge0$, the map sending $i\in\{0,\ldots,n-1\}$ to its right copy when $i<r$ and its left copy otherwise is injective.
\end{lemma}

\begin{lemma}[Initial hybrid equals the left family]
\lean{BooleanAnalysis.GaussianSpace.Invariance.hybridInput_zero}
\leanok
\uses{BooleanAnalysis.GaussianSpace.Invariance.hybridIndex, BooleanAnalysis.GaussianSpace.Invariance.hybridInput}
\difficulty{3}
Let $Z_{L,i},Z_{R,i}\colon\Omega\to\mathbb R$ be two families indexed by $0\le i<n$, and define $H^{(r)}_i=Z_{R,i}$ if $i<r$ and $Z_{L,i}$ otherwise. Then $H^{(0)}_i=Z_{L,i}$ for every $i$.
\end{lemma}

\begin{lemma}[Final hybrid equals the right family]
\lean{BooleanAnalysis.GaussianSpace.Invariance.hybridInput_all}
\leanok
\uses{BooleanAnalysis.GaussianSpace.Invariance.hybridIndex, BooleanAnalysis.GaussianSpace.Invariance.hybridInput}
\difficulty{3}
Let $Z_{L,i},Z_{R,i}\colon\Omega\to\mathbb R$ be two families indexed by $0\le i<n$, and define $H^{(r)}_i=Z_{R,i}$ if $i<r$ and $Z_{L,i}$ otherwise. Then $H^{(n)}_i=Z_{R,i}$ for every $i$, including the empty family when $n=0$.
\end{lemma}

\begin{lemma}[Consecutive hybrids agree away from the replaced input]
\lean{BooleanAnalysis.GaussianSpace.Invariance.hybridInput_succ_off}
\leanok
\uses{BooleanAnalysis.GaussianSpace.Invariance.hybridIndex, BooleanAnalysis.GaussianSpace.Invariance.hybridInput}
\difficulty{3}
Let $Z_{L,j},Z_{R,j}\colon\Omega\to\mathbb R$ be two families indexed by $0\le j<n$, and put $H^{(r)}_j=Z_{R,j}$ if $j<r$ and $Z_{L,j}$ otherwise. If $i,j\in\{0,\ldots,n-1\}$ and $j\ne i$, then $H^{(i)}_j=H^{(i+1)}_j$ as functions on $\Omega$.
\end{lemma}

\begin{theorem}[Invariance for two mutually independent input copies]
\lean{BooleanAnalysis.GaussianSpace.Invariance.invariance_of_independent_pair}
\leanok
\uses{BooleanAnalysis.GaussianSpace.HasInvarianceMoments, BooleanAnalysis.GaussianSpace.Invariance.hybridIndex, BooleanAnalysis.GaussianSpace.Invariance.hybridIndex_injective, BooleanAnalysis.GaussianSpace.Invariance.hybridInput, BooleanAnalysis.GaussianSpace.Invariance.hybridInput_all, BooleanAnalysis.GaussianSpace.Invariance.hybridInput_succ_off, BooleanAnalysis.GaussianSpace.Invariance.hybridInput_zero, BooleanAnalysis.GaussianSpace.single_coordinate_replacement_bound, BooleanAnalysis.Hypercontractivity.MultilinearMoments.eval_supported, BooleanAnalysis.Hypercontractivity.evalMultilinearRV, BooleanAnalysis.ThresholdFunctions.MultilinearPolynomial, BooleanAnalysis.ThresholdFunctions.MultilinearPolynomial.HasDegreeAtMost, BooleanAnalysis.ThresholdFunctions.MultilinearPolynomial.eval, BooleanAnalysis.influence}
\difficulty{5}
Let $F(x)=\sum_{S\subseteq\{1,\ldots,n\}}a_S\prod_{i\in S}x_i$ be a real multilinear polynomial of degree at most $k$; let $I_i(F)=\sum_{S\ni i}a_S^2$. Let $(Z_{L,i},Z_{R,i})_{i=1}^n$ be a mutually independent family of $2n$ real inputs on a probability space, each measurable up to a null set with finite fourth absolute moment, first three moments $0,1,0$, and fourth moment at most $9$. If $\psi\in C^4(\mathbb R)$ and $C\ge0$ bounds $|\psi^{(4)}|$ everywhere, then $|\mathbb E\psi(F(Z_L))-\mathbb E\psi(F(Z_R))|\le(C/12)9^k\sum_i I_i(F)^2$.
\end{theorem}

\begin{theorem}[Joint laws from independent coordinate laws]
\lean{BooleanAnalysis.GaussianSpace.Invariance.identDistrib_of_independent_coordinates}
\leanok
\difficulty{4}
Let $I$ be a finite index set and let $(X_i)_{i\in I}$ and $(Y_i)_{i\in I}$ be real input families on probability spaces $(\Omega,\mu)$ and $(\Xi,\nu)$, respectively. If each family is independent and $X_i$ and $Y_i$ are measurable up to null sets with the same distribution for every $i$, then the random vectors $(X_i)_{i\in I}$ and $(Y_i)_{i\in I}$ have the same distribution.
\end{theorem}
